% ProMarket article — "Same Storm, Different Boats"
% Lodefalk, Löthman, Koch, Engberg (2026)
% For Overleaf editing; submit to ProMarket as plain text/Word.
% American English throughout (ProMarket, Chicago Booth).
% References appear as [Author YEAR] in text — replace with hyperlinks on submission.
% Target: 1,000–1,500 words (current ~1,260).
% Status: DRAFT — requires author rewriting before submission.
%         ProMarket AI policy prohibits substantially AI-written work.
%         This draft is a skeleton for Magnus and Michael Koch to rewrite.

\documentclass[12pt]{article}

\usepackage[utf8]{inputenc}
\usepackage[T1]{fontenc}
\usepackage[margin=1.15in]{geometry}
\usepackage{setspace}
\usepackage{parskip}
\usepackage{microtype}
\usepackage[hidelinks,
  colorlinks=true,
  linkcolor=blue,
  urlcolor=blue,
  citecolor=blue]{hyperref}
\usepackage{booktabs}

\onehalfspacing

\begin{document}

% ─────────────────────────────────────────────────────────────────
% TITLE BLOCK
% ─────────────────────────────────────────────────────────────────

% TITLE OPTIONS — submit 2–3 to ProMarket; editors decide.
% (1) Generative AI Is Changing Who Gets Hired, Not How Many
% (2) The Stock Market Divergence Is Not an AI Story. The Hiring Freeze Is.
% (3) AI Has a Quiet Labor Market Effect. It Is Showing Up in Entry-Level Hiring.

\title{\textbf{Generative AI Is Changing Who Gets Hired, Not How Many}
\\[4pt]
\normalsize\textit{[Working title — ProMarket editors will finalise]}}

\author{
  Magnus Lodefalk\thanks{Örebro University and RATIO Institute.
  \texttt{magnus.lodefalk@oru.se}} \quad
  Lydia Löthman\thanks{Örebro University.} \quad
  Michael Koch\thanks{Aarhus University and Kiel Institute for the World Economy.} \quad
  Erik Engberg\thanks{Örebro University and RATIO Institute.}
}

\date{May 2026 \\ \medskip
\small Draft. Working paper available at
\href{https://www.ai-econlab.com/papers}{ai-econlab.com/papers}.}

\maketitle
\thispagestyle{empty}

\bigskip
\noindent\rule{\textwidth}{0.4pt}
\bigskip

% ─────────────────────────────────────────────────────────────────
% ARTICLE BODY
% ─────────────────────────────────────────────────────────────────

Since late 2022, stock prices have risen sharply while job openings have fallen almost as steeply.
The divergence has become one of the most cited charts in discussions of artificial intelligence and
the labor market, and for many observers the interpretation is straightforward: generative AI is
substituting for workers, depressing labor demand while lifting corporate profits. Our research,
using full-population register data for Sweden and 4.6 million job advertisements, finds that this
reading of the aggregate picture is not supported by the evidence. The underlying concern, however,
is well founded. The effect of generative AI operates below the surface of standard statistics,
through a selective recomposition of who gets hired.

\section*{The aggregate decline is a monetary policy story}

The key is timing. Sweden's central bank, the Riksbank, raised interest rates for the first time
in April 2022, seven months before ChatGPT launched in November 2022. If monetary tightening is the
primary driver of the posting decline, it should begin with the rate hike. If generative AI is the
driver, it should not appear before late 2022 and should be concentrated in high-exposure occupations.

Using the full population of advertisements published on Platsbanken, Sweden's public employment
service vacancy portal, matched to an occupation-level measure of generative-AI exposure
(the \href{https://docs.iza.org/dp16717.pdf}{DAIOE index}, Engberg et al.\ 2024), we find that
postings across all AI-exposure groups declined together from April 2022 onwards. The post-rate-hike
interaction is estimated at $-0.127$ ($p < 0.01$); the post-ChatGPT interaction at $-0.062$ and not
statistically distinguishable from zero ($p = 0.11$). AI-exposed occupations are also uncorrelated
with interest-rate sensitivity, consistent with the two channels being distinct.
\href{https://www.derekthompson.org/p/is-this-the-new-scariest-chart-in}{Thompson (2025)} % VERIFY URL
has argued similarly for the United States: the posting decline tracks the Federal Reserve's
tightening cycle, not the ChatGPT launch.

\section*{Within employers, the age gradient is large and accelerating}

The aggregate picture conceals a different pattern at the level of individual employers.
Using monthly employer-declaration data (the AGI register) covering the full Swedish working
population, linked to four-digit occupation codes and age from the LISA longitudinal register,
we estimate an employer-level difference-in-differences design
(following \href{https://digitaleconomy.stanford.edu/wp-content/uploads/2025/11/CanariesintheCoalMine_Nov25.pdf}{Brynjolfsson et al.\ 2025}). % VERIFY URL
Employer-by-quartile and employer-by-month fixed effects absorb time-invariant compositional
differences and all firm-level time-varying shocks, so identification comes from within-employer
recomposition across AI-exposure quartiles over time.

Employment of workers aged 22 to 25 in the top quartile of AI-exposed occupations falls 5.5~percent
below less-exposed occupations within the same employers by the first half of 2025
(95~percent confidence interval: $-5.8$ to $-4.9$~percent). This is our headline event-study
estimate. Poisson pseudo-maximum-likelihood estimates, which better accommodate the zero-heavy
distribution of employment counts in employer-occupation cells, yield a decline of 16~percent across
four consistent specifications (confidence interval: $-18$ to $-14$~percent). The two estimates are
consistent: the gap in magnitude reflects the scale difference between the log-linear and Poisson
specifications, not a genuine difference in economic magnitude. Workers aged 31 to 49 are essentially
unaffected. Workers over 50 show a small positive gain, though this finding is sensitive to the choice of AI exposure measure.

The adjustment loads almost entirely on the hiring margin. The hires coefficient for ages 22 to 25 is
$-0.0051$ ($p < 0.001$); the separations coefficient is one-quarter as large
($-0.0012$, $p < 0.01$). Sweden's employment-protection legislation, based on a last-in-first-out
principle, constrains dismissal of incumbent workers. High-AI-exposure employers are therefore not
laying off junior staff. Rather, they are not replacing departing workers and not taking on new
cohorts. \href{https://papers.ssrn.com/sol3/papers.cfm?abstract_id=5425555}{Hosseini and Lichtinger (2025)}, % VERIFY SSRN ID
analyzing r\'esum\'e and job-posting data for roughly 65~million workers across more than 280{,}000 U.S.\ firms, find an approximately 9~percent junior employment decline in generative-AI-adopting firms driven by a comparable hiring
slowdown. The hiring-freeze pattern appears consistent across institutional settings.

\section*{Young women face a larger adjustment}

The most under-examined finding concerns young women. Women aged 22 to 25 in high-AI-exposure
occupations have experienced a difference-in-differences employment decline of $-0.016$
($p < 0.001$), compared with $-0.007$ ($p < 0.01$) for men of the same age. The gender gap is
itself statistically significant (triple-difference $p < 0.01$). Roughly two-fifths of this
disparity is accounted for by occupational composition: women in Sweden are concentrated in
administrative, payroll, and customer-service roles near the top of the AI-exposure distribution.
Payroll administrators are 82~percent female and sit at the 96th percentile of the DAIOE index;
customer service agents and receptionists are 66~percent female.

The remaining three-fifths of the gender gap is within-occupation: young women are losing ground
relative to young men in the same AI-exposed jobs. The mechanism is not identified in our data,
but the within-occupation component implies that occupational sorting alone does not account for
the disparity.

This employment finding is consistent with simulation evidence for Sweden from
\href{https://www.ifn.se/en/publications/working-papers/2025/1534/}{Gardberg, Heyman, Olsson and T\aa{}g (IFN~2025)}, % VERIFY URL
who show that pre-AI patterns of gender-based occupational sorting can widen the post-AI gender wage gap.
Our data suggest that employment consequences are already materializing ahead of the wage effects.
The \href{https://www.ilo.org/resource/news/new-ilo-data-confirm-women-face-higher-workplace-risks-generative-ai-men}{International Labour Organization (2026)} % verified live, published 5 March 2026
has estimated that female-dominated occupations globally face roughly twice the generative-AI
exposure risk of male-dominated ones. In Sweden, that elevated risk is now reflected in
employer-level hiring decisions.

\section*{Why outcomes differ across countries}

Three studies, three countries, three outcomes.
\href{https://digitaleconomy.stanford.edu/wp-content/uploads/2025/11/CanariesintheCoalMine_Nov25.pdf}{Brynjolfsson, Chandar and Chen (Stanford Digital Economy Lab 2025)} % VERIFY URL
find a 16~percent Poisson-estimated decline in employment for young U.S.\ workers in AI-exposed
occupations. Our Swedish Poisson estimates cluster in the same range.
\href{https://www.etla.fi/en/publications/working-papers-en/ai-has-not-impacted-the-youth-labor-market-in-finland/}{Kauhanen and Rouvinen (ETLA~2026)}, % VERIFY URL
however, find no effect in Finland using a comparable research design on Finnish population data.
We test directly whether this divergence is methodological: reweighting our Swedish sample to match
Finland's industry-occupation composition and applying Finland's preferred exposure measure leaves
our estimated decline at approximately 15~percent, nearly unchanged. The divergence is not an
artifact of study design.

What explains it remains uncertain. Candidate factors include differences in the intensity of AI
adoption, in the occupational structure of youth employment, and in the vocational-training pathways
that channel young Finns into less-exposed fields. The cross-country heterogeneity is itself
substantive: institutional context appears to shape the distributional consequences of the same
underlying technology.

\section*{A skill-formation concern}

Job counts alone understate what is at stake. Entry-level positions in AI-exposed occupations are
not merely employment. They are the mechanism through which workers accumulate applied skills,
professional networks, and the tacit knowledge that constitutes mid-career capability. An employer
that stops taking on new cohorts today is not only reducing near-term costs. It is also failing to
develop the experienced workers it will need a decade from now.

Brynjolfsson et al.\ (2025) describe young workers as the ``canaries in the coal mine,'' early
signals of a broader adjustment. If the entry-level hiring freeze persists, what follows is not
merely a transient employment dip but a generational skill-formation gap. That gap will appear in
career trajectories and aggregate productivity well before it registers in headline unemployment
figures.

\section*{Implications}

Three practical implications follow. First, monitoring: aggregate unemployment and job-posting
indices will not detect this adjustment early. Employer-level data at monthly frequency, tracking
hiring composition by age and occupation, is what makes it visible. Countries with
employer-declaration registers of the kind Sweden maintains hold a significant informational
advantage. Second, education and training: if entry-level positions in AI-exposed fields contract,
universities and vocational programs face a mandate to restructure applied learning so as to provide
the experiential component that employment has historically supplied. Third, coverage gaps: recent
graduates who cannot find work and have not yet accrued eligibility for unemployment-insurance
benefits are invisible to the social insurance system. Monitoring whether this population is growing
disproportionately in AI-exposed fields warrants attention.

\bigskip
\noindent\rule{\textwidth}{0.4pt}

\noindent\textbf{Disclosure.} Funded by the Torsten S\"oderberg Foundation (grants E46/21, ET3/23)
and WASP-HS (grant 805). Employment data from Statistics Sweden via the MONA platform under ethical
approvals detailed in the \href{https://www.ai-econlab.com/papers}{working paper}. No conflicts of
interest.

% ─────────────────────────────────────────────────────────────────
% FIGURES (send as separate image files to ProMarket)
% ─────────────────────────────────────────────────────────────────
%
% Figure 1: figures/fig1_scary_chart.png
%   Swedish stock market vs job postings by AI-exposure quartile, 2020–2026.
%   Caption: "Sweden's stock market and job postings, 2020–2026. Postings fall across all
%   AI-exposure groups from the Riksbank rate hike (April 2022), not the ChatGPT launch
%   (November 2022)."
%
% Figure 2: figures/fig3_event_study_main.pdf
%   Event study by age group, employer-level, 2019–2025.
%   Caption: "Within-employer employment gap by age group in AI-exposed vs less-exposed
%   occupations, 2019–2025. Young workers (22–25) diverge after ChatGPT; workers over 50
%   are stable. Reference period: first half of 2022."

% ─────────────────────────────────────────────────────────────────
% HYPERLINK CHECKLIST (verify before submission)
% ─────────────────────────────────────────────────────────────────
%
% Thompson 2025 scary chart:   https://www.derekthompson.org/p/is-this-the-new-scariest-chart-in
% Brynjolfsson et al. 2025:    https://digitaleconomy.stanford.edu/wp-content/uploads/2025/11/CanariesintheCoalMine_Nov25.pdf
% Hosseini & Lichtinger 2025:  https://papers.ssrn.com/sol3/papers.cfm?abstract_id=5425555
% Gardberg et al. IFN 2025:    https://www.ifn.se/en/publications/working-papers/2025/1534/
% Kauhanen & Rouvinen 2026:    https://www.etla.fi/en/publications/working-papers-en/ai-has-not-impacted-the-youth-labor-market-in-finland/
% ILO gender AI risk 2026:     https://www.ilo.org/resource/news/new-ilo-data-confirm-women-face-higher-workplace-risks-generative-ai-men  [verified live, published 5 Mar 2026]
% Engberg et al. DAIOE:        https://docs.iza.org/dp16717.pdf
% Working paper:               https://www.ai-econlab.com/papers

\end{document}
